\begin{casebox}
\textbf{知识点小案例：reduce 不能靠扫描临时目录发现 shuffle。}
特例是 map attempt 0 写出 \code{block.tmp} 后 reply 丢失，attempt 1 重试并成功提交。目录里可能同时存在两个物理 block，其中一个是未提交或迟到产物。若 reduce 直接扫描目录，它无法知道哪个 block 属于被接受的 attempt。manifest 记录 logical partition、attempt、checksum、size、commit time 和状态，reduce 只读取 committed entry。由此推出规律：数据面文件存在不等于控制面事实成立。工程上所有 shuffle、checkpoint 和推理模型 artifact 都应通过 manifest 建立可恢复提交边界。
\end{casebox}
